\documentclass{article}
\usepackage{amsmath,amssymb,amsthm}

\newtheorem{lemma}{Lemma}
\newtheorem{remark}{Remark}

\begin{document}

\section*{Status note}

The benchmark record \texttt{rm-02895} contains \emph{no} problem statement:
the issue body is empty boilerplate (\textit{Area: (unspecified)},
\textit{Problem author: (unknown)}), there is no associated
\texttt{problem.tex}, and the registered problem set is \texttt{q1}--\texttt{q10}
only. Consequently there is no author-intended proposition to prove, and the
Faithfulness requirement forbids substituting a different statement for a
missing one. The argument below is therefore presented \emph{explicitly as a
placeholder}: it is the standard ``Universal Chord Theorem'' proposed earlier
in the thread, written out in full and rigorous form so that the thread holds a
complete, correct artifact. It is \emph{not} a recovery of the (nonexistent)
intended problem.

\section*{Placeholder Lemma (Universal Chord Theorem)}

\begin{lemma}
Let $f:[0,1]\to\mathbb{R}$ be continuous with $f(0)=f(1)$, and let
$n\in\mathbb{N}=\{1,2,3,\dots\}$. Then there exists
$x\in\bigl[0,\,1-\tfrac1n\bigr]$ such that
\[
  f\!\left(x+\tfrac1n\right)=f(x).
\]
\end{lemma}

\begin{proof}
Define
\[
  g:\Bigl[0,\,1-\tfrac1n\Bigr]\to\mathbb{R},\qquad
  g(x)=f\!\left(x+\tfrac1n\right)-f(x).
\]
Since $f$ is continuous on $[0,1]$ and $x\mapsto x+\tfrac1n$ is continuous,
$g$ is continuous on $\bigl[0,1-\tfrac1n\bigr]$.

\medskip
\noindent\textbf{The grid points and the telescoping identity.}
For $k=0,1,\dots,n-1$ the point $\tfrac{k}{n}$ lies in
$\bigl[0,1-\tfrac1n\bigr]$ (indeed $0\le \tfrac kn\le \tfrac{n-1}{n}=1-\tfrac1n$),
so $g(\tfrac kn)$ is defined, and moreover both
$\tfrac kn$ and $\tfrac{k+1}{n}=\tfrac kn+\tfrac1n$ lie in $[0,1]$, the domain
of $f$. We compute the telescoping sum
\[
  \sum_{k=0}^{n-1} g\!\left(\tfrac kn\right)
   =\sum_{k=0}^{n-1}\Bigl(f\!\left(\tfrac{k+1}{n}\right)-f\!\left(\tfrac kn\right)\Bigr)
   = f\!\left(\tfrac nn\right)-f\!\left(\tfrac 0n\right)
   = f(1)-f(0) = 0. \tag{$\ast$}
\]

\medskip
\noindent\textbf{Sign argument.}
We distinguish two cases.

\emph{Case 1: some $g(\tfrac kn)=0$.} Then $x=\tfrac kn$ satisfies the
conclusion and we are done.

\emph{Case 2: $g(\tfrac kn)\neq 0$ for every $k\in\{0,\dots,n-1\}$.}
If all these $n$ nonzero numbers had the same sign, their sum would be
strictly positive (if all $>0$) or strictly negative (if all $<0$),
contradicting $(\ast)$. (Here we use $n\ge 1$, so there is at least one term;
in fact for $n=1$ identity $(\ast)$ already reads $g(0)=0$, placing us in
Case~1.) Hence not all of $g(\tfrac0n),\dots,g(\tfrac{n-1}{n})$ share a common
sign, so there exist consecutive indices $k$ and $k+1$ (with
$0\le k\le n-2$) for which
\[
  g\!\left(\tfrac kn\right)\quad\text{and}\quad g\!\left(\tfrac{k+1}{n}\right)
\]
have opposite (strictly nonzero) signs, i.e.
$g(\tfrac kn)\,g(\tfrac{k+1}{n})<0$.

Indeed: list the signs $\varepsilon_k=\operatorname{sgn} g(\tfrac kn)\in\{+1,-1\}$
for $k=0,\dots,n-1$. They are not all equal, so as $k$ increases from $0$ to
$n-1$ the sign changes at least once between two consecutive indices; pick such
a pair.

\medskip
\noindent\textbf{Intermediate Value Theorem.}
Both grid points $\tfrac kn$ and $\tfrac{k+1}{n}$ lie in
$\bigl[0,1-\tfrac1n\bigr]$ (since $0\le k\le k+1\le n-1$), so the closed
interval $\bigl[\tfrac kn,\tfrac{k+1}{n}\bigr]$ is contained in the domain of
$g$, on which $g$ is continuous. Because $g(\tfrac kn)$ and $g(\tfrac{k+1}{n})$
have opposite signs, the Intermediate Value Theorem yields a point
$x\in\bigl(\tfrac kn,\tfrac{k+1}{n}\bigr)\subseteq\bigl[0,1-\tfrac1n\bigr]$ with
$g(x)=0$, that is, $f(x+\tfrac1n)=f(x)$.

In both cases we have produced $x\in[0,1-\tfrac1n]$ with $f(x+\tfrac1n)=f(x)$,
proving the lemma.
\end{proof}

\begin{remark}[Sharpness in the step size]
The conclusion is special to step sizes of the form $h=\tfrac1n$. For a step
$h\in(0,1)$ that is \emph{not} the reciprocal of an integer, it can fail. A
standard counterexample: fix such an $h$ and put
\[
  f(x)=\sin^2\!\Bigl(\tfrac{\pi x}{h}\Bigr)-x\,\sin^2\!\Bigl(\tfrac{\pi}{h}\Bigr).
\]
Then $f$ is continuous, $f(0)=0$ and $f(1)=\sin^2(\pi/h)-\sin^2(\pi/h)=0$, yet a
direct computation gives, for all $x\in[0,1-h]$,
\[
  f(x+h)-f(x)=-h\,\sin^2\!\Bigl(\tfrac{\pi}{h}\Bigr)<0,
\]
because $\sin^2(\pi/h)\neq0$ when $1/h\notin\mathbb{Z}$. Hence no chord of
horizontal length $h$ exists. (Numerically verified for $h=0.4$:
$f(x+h)-f(x)\equiv-0.4$.)
\end{remark}

\begin{remark}[Scope]
This document resolves only the meta-blocker (no statement was present) by
recording a complete, verified placeholder argument. The underlying tracker
entry \texttt{rm-02895} still lacks an actual problem statement; a genuine
resolution requires the intended proposition to be supplied.
\end{remark}

\end{document}
